\documentclass[11pt]{article}
\usepackage[margin=1in]{geometry}
\usepackage{enumitem}
% =============================================================================
% Preambles/header.tex — shared LaTeX/Beamer preamble for this project
%
% Use from a Beamer lecture by prepending:
%     \input{header}
% and compiling with TEXINPUTS=../Preambles:$TEXINPUTS (the /compile-latex
% skill does this automatically).
%
% PALETTE CONTRACT. Color names defined here MUST match the names in
% Quarto/theme-template.scss so Beamer and Quarto renderings use the same
% palette. The scripts/check-palette-sync.sh script enforces this.
%
% To customize: edit the HEX values below AND the matching $variable in
% Quarto/theme-template.scss, then run scripts/check-palette-sync.sh.
% =============================================================================

% -----------------------------------------------------------------------------
% Engine detection + encoding
%
% Our /compile-latex skill uses XeLaTeX, where inputenc is unnecessary and
% can produce encoding warnings. Only load inputenc under pdfTeX.
% -----------------------------------------------------------------------------
\usepackage{iftex}
\ifPDFTeX
  \usepackage[utf8]{inputenc}
\fi

\usepackage{amsmath, amssymb, amsthm}
\usepackage{booktabs}
\usepackage{graphicx}

% -----------------------------------------------------------------------------
% PALETTE — mirrors Quarto/theme-template.scss variable names.
% Load xcolor and define colors BEFORE anything that references them
% (notably hyperref, hypersetup, and the Beamer color assignments).
% -----------------------------------------------------------------------------
\usepackage{xcolor}

\definecolor{primary-blue}{HTML}{012169}      % headings, accents
\definecolor{primary-gold}{HTML}{B9975B}      % emphasis, borders
\definecolor{highlight-yellow}{HTML}{F2A900}  % markers, alerts
\definecolor{light-bg}{HTML}{E8EDF5}          % title / section backgrounds
\definecolor{jet}{HTML}{1A1A1A}               % body text

% Semantic colors (used in diagrams and annotations)
\definecolor{positive}{HTML}{15803D}          % good / observed / identified
\definecolor{negative}{HTML}{B91C1C}          % bad / problematic / bias
\definecolor{neutral}{HTML}{525252}           % reference / context

% Highlight accents (match .hi-* classes in the SCSS)
\definecolor{hi-slate}{HTML}{314F4F}
\definecolor{hi-green}{HTML}{2E8B57}
\definecolor{hi-red}{HTML}{C41E3A}

% Semantic aliases so skill templates and snippets can write
% \textcolor{accent}{...} without hard-coding "primary-blue".
\colorlet{accent}{primary-blue}
\colorlet{accent2}{primary-gold}

% -----------------------------------------------------------------------------
% Hyperref — loaded AFTER xcolor + palette so link colors resolve.
% -----------------------------------------------------------------------------
\usepackage{hyperref}
\hypersetup{colorlinks=true, linkcolor=primary-blue, urlcolor=primary-blue,
            citecolor=primary-blue, pdfborder={0 0 0}}

% -----------------------------------------------------------------------------
% Course code — set once per deck (after \input{header}, before \begin{document})
% via \coursecode{CS401}. Shown in the footer on every slide so a deck's course
% is identifiable without opening the title page. Empty by default.
% -----------------------------------------------------------------------------
\makeatletter
\newcommand{\coursecode}[1]{\gdef\@coursecodetext{#1}}
\gdef\@coursecodetext{}
\makeatother

% -----------------------------------------------------------------------------
% Beamer theme assignments (only apply under Beamer).
%
% We detect Beamer via \@ifclassloaded{beamer}, which is the documented,
% non-internal way. This must be wrapped in \makeatletter/\makeatother
% because \@ifclassloaded contains an @-catcode character.
% -----------------------------------------------------------------------------
\makeatletter
\@ifclassloaded{beamer}{%
  \usefonttheme{professionalfonts}
  \setbeamercolor{structure}{fg=primary-blue}
  \setbeamercolor{frametitle}{fg=primary-blue}
  \setbeamercolor{title}{fg=primary-blue}
  \setbeamercolor{subtitle}{fg=primary-gold}
  \setbeamercolor{author}{fg=primary-blue}
  \setbeamercolor{institute}{fg=neutral}
  \setbeamercolor{date}{fg=primary-gold}
  \setbeamercolor{itemize item}{fg=primary-blue}
  \setbeamercolor{itemize subitem}{fg=primary-gold}
  \setbeamercolor{alerted text}{fg=primary-gold}
  \setbeamercolor{block title}{fg=white, bg=primary-blue}
  \setbeamercolor{block body}{bg=light-bg}
  % Semantic block variants: block=definition/concept (blue), exampleblock=
  % worked example (green/positive), alertblock=key takeaway or caution (gold).
  % Keep to <=2 colored boxes per slide across ALL three variants (INV-7).
  \setbeamercolor{block title example}{fg=white, bg=positive}
  \setbeamercolor{block body example}{bg=light-bg}
  \setbeamercolor{block title alerted}{fg=white, bg=primary-gold}
  \setbeamercolor{block body alerted}{bg=light-bg}

  % Minimal footer: course code (left) + page X / N (right), no navigation clutter
  \setbeamertemplate{navigation symbols}{}
  \setbeamertemplate{footline}{%
    \color{neutral}\footnotesize\hspace{0.7em}\@coursecodetext\hfill
    \insertframenumber/\inserttotalframenumber\hspace{0.7em}\vspace{0.3em}}
}{}
\makeatother

% -----------------------------------------------------------------------------
% TikZ setup — libraries the snippet gallery uses
% -----------------------------------------------------------------------------
\usepackage{tikz}
\usetikzlibrary{arrows.meta, positioning, calc, decorations.pathreplacing,
                fit, shapes.geometric, backgrounds}

% Shared TikZ styles — snippets and hand-written diagrams can pick these up
% via `\begin{tikzpicture}[every node/.style={...}]` or by naming specific
% styles (e.g., `\node[dag-node] (X) at (0,0) {X};`).
\tikzset{
  % Node styles for causal DAGs and similar
  dag-node/.style = {
    draw=primary-blue, fill=light-bg, rounded corners=2pt,
    minimum width=1.2cm, minimum height=0.9cm, align=center, font=\small
  },
  decision-node/.style = {
    draw=primary-gold, fill=white, diamond, aspect=1.8,
    minimum width=1.8cm, minimum height=1.2cm, align=center, font=\small,
    inner sep=1pt
  },
  % Edge styles — observed vs counterfactual
  observed-edge/.style = {-{Latex[length=2.5mm]}, thick, draw=jet},
  counterfactual-edge/.style = {-{Latex[length=2.5mm]}, thick, draw=neutral, dashed},
  confound-edge/.style = {-{Latex[length=2.5mm]}, thick, draw=negative, dashed},
  % Data-point styles for plots
  observed-dot/.style = {fill=primary-blue, draw=primary-blue, circle, inner sep=1.2pt},
  counterfactual-dot/.style = {fill=white, draw=neutral, circle, inner sep=1.2pt}
}

% -----------------------------------------------------------------------------
% Convenience macros
% -----------------------------------------------------------------------------
\newcommand{\muted}[1]{\textcolor{neutral}{#1}}
\newcommand{\key}[1]{\textcolor{primary-gold}{\textbf{#1}}}
\newcommand{\good}[1]{\textcolor{positive}{#1}}
\newcommand{\bad}[1]{\textcolor{negative}{#1}}

% Standout frame macro: full-bleed dark background for section transitions.
% Beamer-only (uses \begin{frame}/\setbeamercolor/\usebeamerfont) -- do not
% call from an article-class file (e.g. Notes/<CODE>/*-notes.tex). \muted,
% \key, \good, \bad, and \coursecode above are plain xcolor/gdef and are
% safe to use from either class; verified when Notes/article-class support
% was added.
% Expand: \transitionslide{Your transition title}
\newcommand{\transitionslide}[1]{%
  \begingroup
  \setbeamercolor{background canvas}{bg=primary-blue}
  \begin{frame}[plain]
    \centering
    \vfill
    {\usebeamerfont{title}\color{white}\Large #1\par}
    \vfill
  \end{frame}
  \endgroup
}

% Section divider frame (Beamer-only), an alternative to \transitionslide with a
% darker two-line style: near-black (jet) background, a small white label line
% above a larger gold title, and a thin gold rule along the bottom edge.
% Expand: \sectiondivider{Section 2}{Pointers}
\newcommand{\sectiondivider}[2]{%
  \begingroup
  \setbeamercolor{background canvas}{bg=jet}
  \begin{frame}[plain]
    \vspace*{3.0cm}
    {\color{white}\ttfamily\large #1\par}
    \vspace{0.5cm}
    {\color{primary-gold}\LARGE #2\par}
    \begin{tikzpicture}[remember picture, overlay]
      \draw[primary-gold, line width=2.5pt]
        ($(current page.south west)+(0,0.1)$) -- ($(current page.south east)+(0,0.1)$);
    \end{tikzpicture}
  \end{frame}
  \endgroup
}

% -----------------------------------------------------------------------------
% Channel watermark — "Concepts That Click" (YouTube).
%
% Article-class files (CompetitiveExam Q/A sets, Notes, Assignments) get a
% light diagonal draft watermark on every page plus a centered footer naming
% the channel. Beamer decks get a subtle TikZ background watermark instead
% (the footline already carries the course code). Centralized here so every
% MCQ PDF is branded without per-file edits.
% -----------------------------------------------------------------------------
\makeatletter
\@ifclassloaded{article}{%
  \usepackage{draftwatermark}
  \SetWatermarkText{Concepts That Click}
  \SetWatermarkScale{1}
  \SetWatermarkFontSize{2cm}
  \SetWatermarkLightness{0.86}
  \SetWatermarkAngle{45}
  \usepackage{fancyhdr}
  \pagestyle{fancy}
  \fancyhf{}
  \fancyfoot[C]{\footnotesize\textcolor{neutral}{Concepts That Click~|~YouTube: \href{https://www.youtube.com/@concepts.thatclick}{@concepts.thatclick}}}
  \renewcommand{\headrulewidth}{0pt}
  \renewcommand{\footrulewidth}{0pt}
  \fancypagestyle{plain}{%
    \fancyhf{}%
    \fancyfoot[C]{\footnotesize\textcolor{neutral}{Concepts That Click~|~YouTube: \href{https://www.youtube.com/@concepts.thatclick}{@concepts.thatclick}}}%
    \renewcommand{\headrulewidth}{0pt}%
    \renewcommand{\footrulewidth}{0pt}%
  }
}{}
\@ifclassloaded{beamer}{%
  \setbeamertemplate{background}{%
    \begin{tikzpicture}[remember picture, overlay]
      \node[rotate=45, opacity=0.055, align=center]
        at (current page.center)
        {\Huge\textcolor{neutral}{Concepts That Click}};
    \end{tikzpicture}%
  }
}{}
\makeatother 
\coursecode{JKSSB}
\title{Lesson 05 Practice: CPU, ALU, Control Unit \& Registers}
\author{JKSSB Computer Awareness}
\date{\today}
\begin{document}
\maketitle
\noindent\textbf{Provenance:} 20 original JKSSB-pattern questions
(\textit{[Original]}) written for this lesson; none reproduces an official
paper item.

\begin{enumerate}[label=\textbf{Q\arabic*.}, leftmargin=*]
\item \textit{[Original]} Expand ALU.
  \begin{enumerate}[label=\Alph*.]
    \item Arithmetic Logic Unit
    \item Advanced Logic Unit
    \item Arithmetic Language Unit
    \item Application Logic Unit
  \end{enumerate}
\item \textit{[Original]} Which CPU component performs arithmetic and
  logical operations on data?
  \begin{enumerate}[label=\Alph*.]
    \item Control Unit
    \item Register file
    \item Arithmetic Logic Unit (ALU)
    \item System bus
  \end{enumerate}
\item \textit{[Original]} The main function of the Control Unit is to:
  \begin{enumerate}[label=\Alph*.]
    \item perform addition and subtraction of numbers
    \item direct the sequence of operations and issue control signals
    \item store intermediate arithmetic results
    \item provide permanent secondary storage
  \end{enumerate}
\item \textit{[Original]} Which of the following is \textbf{NOT} a function
  of the Control Unit?
  \begin{enumerate}[label=\Alph*.]
    \item Decoding instructions
    \item Issuing control signals
    \item Directing the flow of data between units
    \item Performing the addition of two numbers
  \end{enumerate}
\item \textit{[Original]} Which register holds the address of the
  \textbf{next} instruction to be fetched?
  \begin{enumerate}[label=\Alph*.]
    \item Instruction Register (IR)
    \item Program Counter (PC)
    \item Memory Data Register (MDR)
    \item Accumulator (ACC)
  \end{enumerate}
\item \textit{[Original]} Which register holds the instruction that is
  currently being decoded and executed?
  \begin{enumerate}[label=\Alph*.]
    \item Program Counter (PC)
    \item Memory Data Register (MDR)
    \item Instruction Register (IR)
    \item Memory Address Register (MAR)
  \end{enumerate}
\item \textit{[Original]} Which register holds the address of the memory
  location to be accessed?
  \begin{enumerate}[label=\Alph*.]
    \item Memory Data Register (MDR)
    \item Memory Address Register (MAR)
    \item Instruction Register (IR)
    \item Accumulator (ACC)
  \end{enumerate}
\item \textit{[Original]} Which register holds the data being transferred to
  or from memory?
  \begin{enumerate}[label=\Alph*.]
    \item Memory Address Register (MAR)
    \item Program Counter (PC)
    \item Memory Data Register (MDR)
    \item Instruction Register (IR)
  \end{enumerate}
\item \textit{[Original]} The accumulator (ACC) is primarily used to hold:
  \begin{enumerate}[label=\Alph*.]
    \item the address of the next instruction
    \item intermediate arithmetic and logic results
    \item the last instruction fetched from memory
    \item a permanent copy of the operating system
  \end{enumerate}
\item \textit{[Original]} Which is the fastest storage in a computer system?
  \begin{enumerate}[label=\Alph*.]
    \item Registers
    \item L1 cache
    \item RAM
    \item Solid-state drive
  \end{enumerate}
\item \textit{[Original]} Where are registers located?
  \begin{enumerate}[label=\Alph*.]
    \item In cache memory
    \item In main memory (RAM)
    \item In secondary storage
    \item Inside the CPU
  \end{enumerate}
\item \textit{[Original]} Clock speed is measured in hertz (Hz). A processor
  rated at 2.4 GHz completes approximately:
  \begin{enumerate}[label=\Alph*.]
    \item 2.4 million cycles per second
    \item 2.4 billion cycles per second
    \item 2.4 thousand cycles per second
    \item 2.4 trillion cycles per second
  \end{enumerate}
\item \textit{[Original]} Which statement about clock cycles is correct?
  \begin{enumerate}[label=\Alph*.]
    \item Every machine instruction completes in exactly one clock cycle.
    \item All instructions require the same number of clock cycles.
    \item Different instructions may take different numbers of clock cycles.
    \item A clock cycle is measured in bytes.
  \end{enumerate}
\item \textit{[Original]} Evaluate the statements:
  \begin{enumerate}[label=\Roman*.]
    \item The ALU works on binary data using logic gates.
    \item The Control Unit directs the sequence of operations and data flow.
    \item The Control Unit performs the arithmetic calculations.
  \end{enumerate}
  \begin{enumerate}[label=\Alph*.]
    \item I only
    \item I and II only
    \item II and III only
    \item I, II, and III
  \end{enumerate}
\item \textit{[Original]} Which abbreviation is correctly matched with its
  full form?
  \begin{enumerate}[label=\Alph*.]
    \item PC --- Program Counter
    \item IR --- Instruction Recorder
    \item MAR --- Memory Arithmetic Register
    \item MDR --- Main Data Register
  \end{enumerate}
\item \textit{[Original]} The Program Counter holds the address of the
  \underline{\hspace{2.2cm}} instruction to be fetched.
  \begin{enumerate}[label=\Alph*.]
    \item next
    \item current
    \item previously executed
    \item data
  \end{enumerate}
\item \textit{[Original]} Registers used for general temporary storage of
  operands and intermediate results, without a fixed dedicated role, are
  called:
  \begin{enumerate}[label=\Alph*.]
    \item status registers
    \item general-purpose registers
    \item memory address registers
    \item instruction registers
  \end{enumerate}
\item \textit{[Original]} The status (flags) register holds:
  \begin{enumerate}[label=\Alph*.]
    \item the address of the memory location to access
    \item the instruction to be decoded
    \item condition flags such as zero, carry, and sign
    \item the data to be written to memory
  \end{enumerate}
\item \textit{[Original]} A quad-core processor means that the processor:
  \begin{enumerate}[label=\Alph*.]
    \item has four bytes of cache memory
    \item supports four monitors at once
    \item runs four operating systems
    \item contains four independent processing cores on a single chip
  \end{enumerate}
\item \textit{[Original]} In the office desktop, which unit fetches each
  instruction, decodes it, and issues control signals to the ALU, registers,
  and memory?
  \begin{enumerate}[label=\Alph*.]
    \item Arithmetic Logic Unit
    \item Control Unit
    \item Hard disk
    \item Monitor
  \end{enumerate}
\end{enumerate}
\end{document}
